%%%%%%%%%%%%%%%%%%%%%%%%%%%%%%%%%%%%%%%%%%%%%%%%%%%%%%%%%%%%%%%%%
% PHIẾU HỌC TẬP: WAVEFORM
%%%%%%%%%%%%%%%%%%%%%%%%%%%%%%%%%%%%%%%%%%%%%%%%%%%%%%%%%%%%%%%%%
\documentclass[12pt,a4paper]{article}

\usepackage[utf8]{inputenc}
\usepackage[T1]{fontenc}
\usepackage[vietnamese]{babel}
\usepackage{amsmath,amssymb}
\usepackage{geometry}
\usepackage{tabularx, array}
\usepackage[svgnames]{xcolor}
\usepackage{tikz}
\usepackage{pgfplots}
\pgfplotsset{compat=1.18}

\geometry{left=2cm,right=2cm,top=2cm,bottom=2cm}
\definecolor{wavepurple}{RGB}{139, 92, 246}

\begin{document}

\begin{center}
    \begin{tikzpicture}
        \node[draw=wavepurple, line width=2pt, rounded corners=10pt, inner sep=10pt] (box) {
            \begin{minipage}{0.9\textwidth}
                \centering
                \bfseries\Large\color{wavepurple} PHÒNG THÍ NGHIỆM ÂM THANH: WAVEFORM
            \end{minipage}
        };
    \end{tikzpicture}
    \vspace{0.5cm}
\end{center}

\noindent \textbf{Nhạc sĩ Toán học (Học sinh):} \dotfill \textbf{Lớp:} \dotfill

\section*{\color{wavepurple} PHẦN 1: QUAN SÁT SÓNG}
Sử dụng App WaveForm để bấm các nốt và điền thông tin:

\begin{center}
\begin{tabular}{|c|c|c|c|}
\hline
\textbf{Nốt nhạc} & \textbf{Tần số $f$ (Hz)} & \textbf{Chu kỳ $T = 1/f$ (s)} & \textbf{Nhận xét Sóng} \\
\hline
C4 (Đồ) & 262 & $\approx 0.00382$ & Sóng thưa \\
\hline
A4 (La) & 440 & & \\
\hline
C5 (Đồ cao) & & & \\
\hline
\end{tabular}
\end{center}

\textit{Nhận xét: Nốt càng cao thì sóng càng} .................

\section*{\color{wavepurple} PHẦN 2: TÍNH TOÁN}
\textbf{Công thức:} $y(t) = A \sin(2\pi f t)$

\textbf{Bài 1:} Cho nốt G4 có tần số $f = 392$ Hz. Tính:
\begin{itemize}
    \item Chu kỳ $T$: $T = \frac{1}{f} = \frac{1}{392} = $ ......... (s)
    \item Trong 1 giây, sóng hoàn thành bao nhiêu dao động? .........
\end{itemize}

\textbf{Bài 2:} Viết phương trình sóng cho nốt E4 ($f=330$ Hz) với biên độ $A=1.5$:
$$ y(t) = ......... \sin(2\pi \cdot ......... \cdot t) $$

\textbf{Bài 3:} Nốt A4 (La chuẩn) có $f=440$ Hz. Nốt A5 (La cao hơn 1 Octave) có tần số bằng bao nhiêu?
$$ f_{A5} = 2 \times f_{A4} = 2 \times 440 = ......... \text{ Hz} $$

\section*{\color{wavepurple} PHẦN 3: VẼ ĐỒ THỊ}
Vẽ phác thảo đồ thị của 2 sóng sin trong khoảng $0 \le t \le 0.01$ (s):
\begin{itemize}
    \item Sóng 1: $y_1 = \sin(2\pi \cdot 100 \cdot t)$ (Tần số thấp)
    \item Sóng 2: $y_2 = \sin(2\pi \cdot 300 \cdot t)$ (Tần số cao)
\end{itemize}

\begin{center}
\begin{tikzpicture}
\begin{axis}[
    width=14cm, height=5cm,
    domain=0:0.01, samples=200,
    xlabel={$t$ (s)}, ylabel={$y$},
    grid=major,
    ymin=-1.5, ymax=1.5,
    xtick={0, 0.002, 0.004, 0.006, 0.008, 0.01}
]
% Students will draw here
\end{axis}
\end{tikzpicture}
\end{center}

\textit{So sánh:} Sóng 2 có số "đỉnh" nhiều hơn sóng 1 trong cùng khoảng thời gian vì ..........................

\end{document}
